\chapter{Instrumentos virtuales modulares}
\section{\obj}
\capacidad
\begin{itemize}
\item Crear instrumentos virtuales modulares para simplificar la creación, mantenimiento, documentación y automatización de los procesos de medición.
\end{itemize}

\section{\mat}
\textbf{A suministrar por la Escuela:}
\begin{itemize}
\item 1 dispositivo de adquisición de datos NI myDAQ
\item 1 termistor TDC310 de \SI{10}{\kilo\ohm} 
\item 1 resistor de \SI{10}{\kilo\ohm}
\item 1 capacitor de \SI{0.1}{\micro\farad}

\end{itemize}
\textbf{A suministrar por el estudiante:}
\begin{itemize}
\item 1 breadboard
\item Cables de interconexión (jumpers) macho-macho
\end{itemize}

\section{\pro}
\begin{enumerate}
\item Realice las conexiones del circuito tal y como se indica en la Figura \ref{fig:L3F1}.
\item Cree un instrumento virtual (\emph{.vi}) y realice lo siguiente:
    \begin{itemize}
        \item A partir de la máquina de estados creada en la Práctica 5, cree tres módulos:
            \begin{itemize}
                \item Un módulo que se encargue de la medición de temperatura, una de sus salidas debe ser un arreglo con las mediciones realizadas.  
                \item Un módulo que se encargue de crear los archivos (\emph{.csv}) y guardar los datos en ellos. 
                \item Un módulo que incluya la maquina de estados completa e internamente contenga los dos módulos anteriores.
                \item Un módulo que configure la adquisición de datos.
            \end{itemize}
    \end{itemize}

\item Incluya en su informe lo siguiente:
    \begin{itemize}
        \item Una captura de pantalla de cada módulo 
        \item Una captura de pantalla del panel frontal
    \end{itemize}
    
\end{enumerate}
